% chapters/12/section-06.tex
\section{Padanan Konsep dengan Pygame}
\label{sec:padanan-unity}

Inilah bagian yang menjelaskan mengapa blok Python tidak menjadi pekerjaan
yang sia-sia. Seluruh konsep yang telah Anda pelajari tetap berlaku; yang
berubah hanyalah nama dan cara memakainya. Tabel~\ref{tab:padanan-lengkap}
memuat padanannya secara lengkap.

\begin{table}[htbp]
\centering
\caption{Padanan lengkap konsep Pygame dan Unity}
\label{tab:padanan-lengkap}
{\setlength{\tabcolsep}{4pt}\footnotesize%
\begin{tabularx}{\textwidth}{@{}L{3.4cm} L{4.0cm} Y@{}}
\toprule
\headrow \textbf{Konsep} & \textbf{Pada Pygame} & \textbf{Pada Unity} \\
\midrule
Siklus permainan & Perulangan yang Anda tulis (Bab~IV) &
Dijalankan mesin; Anda mengisi fungsi \code{Update}. \\
\addlinespace
Selang waktu & \code{delta_seconds} hasil pengukuran \code{clock.tick()} &
\code{Time.deltaTime}. \\
\addlinespace
Keadaan benda & Variabel pada objek Python &
Nilai \code{transform.position}, \code{transform.rotation}, dan peubah pada
skrip. \\
\addlinespace
Gambar benda & \code{sprite.image} &
Komponen \istilah{Sprite Renderer} beserta aset gambarnya. \\
\addlinespace
Persegi pembatas & \code{sprite.rect} &
\istilah{Collider} beserta ukurannya. \\
\addlinespace
Gerak & \code{position += velocity * dt} (Bab~V) &
\code{transform.position += direction * speed * Time.deltaTime}. \\
\addlinespace
Fisika & Perhitungan sendiri (Bab~VI) &
Komponen \istilah{Rigidbody}. \\
\addlinespace
Tabrakan & \code{spritecollide} dan penyelesaian per sumbu (Bab~VII) &
Diperiksa mesin dan dilaporkan sebagai peristiwa \\
\addlinespace
Kelas benda & Kelas turunan \code{pygame.sprite.Sprite} (Bab~VII) &
\istilah{GameObject} beserta komponennya. \\
\addlinespace
Wadah benda & \code{pygame.sprite.Group} &
Daftar pada \istilah{Hierarchy} atau kumpulan pada skrip. \\
\addlinespace
Tingkat permainan & Keadaan yang disusun di dalam kode &
Istilah{Scene}. \\
\addlinespace
Keluaran pesan & \code{print} pada terminal &
\code{Debug.Log} pada panel \istilah{Console}. \\
\bottomrule
\end{tabularx}}
\end{table}

\subsection{Cara Membaca Tabel Itu}

Baris yang paling banyak berubah adalah baris yang berkaitan dengan
\textbf{siapa yang mengerjakan}. Perhatikan tiga contoh.

\begin{enumerate}
  \item Pada baris \istilah{siklus permainan}, pekerjaan yang pada Bab~IV
        menuntut puluhan baris kode kini tidak ditulis sama sekali. Yang perlu
        Anda pahami bukan lagi cara menulisnya, melainkan \textbf{kapan} fungsi
        Anda dipanggil.
  \item Pada baris \istilah{tabrakan}, pemeriksaannya diserahkan kepada mesin.
        Namun penyelesaiannya---yaitu menentukan apa yang terjadi ketika dua
        benda bersentuhan---tetap menjadi tanggung jawab Anda, sebagaimana
        dijelaskan pada Bagian~7.2.
  \item Pada baris \istilah{kelas benda}, perbedaan cara menyusunnya paling
        besar. Objek pada Python disusun dengan mewarisi kelas, sedangkan
        \istilah{GameObject} pada Unity disusun dengan menempelkan komponen.
        Keduanya bertujuan sama, yaitu menyatukan data dan perilaku.
\end{enumerate}

\subsection{Konsep yang Tetap Sama Persis}

Beberapa konsep tidak berubah sama sekali, dan inilah bekal terpenting dari blok
Python.

\begin{itemize}
  \item \textbf{Kecepatan dinyatakan dalam satuan per detik} dan dikalikan
        selang waktu. \code{Time.deltaTime} menggantikan
        \code{delta_seconds}, tetapi perannya sama persis.
  \item \textbf{Arah dinormalkan sebelum dikalikan kecepatan}, agar gerak
        diagonal tidak lebih cepat.
  \item \textbf{Posisi yang sebenarnya disimpan sebagai pecahan}, sedangkan
        yang dipakai memeriksa tabrakan berupa bilangan bulat. Pada Unity,
        keduanya disediakan mesin.
  \item \textbf{Mendeteksi tabrakan dan menyelesaikannya adalah dua pekerjaan
        yang berbeda}. Yang pertama dapat diserahkan kepada mesin; yang kedua
        tidak.
  \item \textbf{Perhitungan yang mahal tetapi tidak mendesak sebaiknya
        disebar}, bukan dijalankan setiap bingkai (Bab~XI).
\end{itemize}

\begin{catatabox}
Kelima butir di atas tidak dapat dipelajari dari panduan
Unity mana pun, sebab panduan itu membahas \textbf{cara memakai} Unity, bukan
\textbf{cara kerja} permainan. Pengetahuan itulah yang Anda peroleh dengan
menulis seluruh mekanismenya sendiri pada blok Python, dan itulah alasan urutan
materi pada mata kuliah ini disusun demikian.
\end{catatabox}
